%------------------------
% Cover Letter Template — Satvik Jain
% Design: Horizontal gradient header band + clean letter body
% Inspired by: Simple Cover Letter (Sid Lacy / Overleaf, CC BY 4.0)
%
% WCAG CONTRAST NOTES (white text on cprimary band):
%   Name at 26pt bold scshape  = WCAG "large text" -> 3:1 minimum required.
%   Contact row at ~9pt small  = standard text     -> 4.5:1 minimum required.
%   Google Blue #4285F4 vs white: ~4.2:1 (passes large text; near-miss small).
%   Gradient darkens 45% toward right edge, improving contrast there.
%   For strict AA on contact row: choose cprimary with luminance below 0.18
%   (example: darken Google Blue to #3A78E0 -> ~4.7:1 with white).
%   Checker: https://webaim.org/resources/contrastchecker/
%
% Rules:  No em-dashes or en-dashes anywhere in body text.
%         No AI-smell openers ("I am writing to express", "I am excited to").
%         Achievement labels appear in csection color; body text stays black.
%         Replace ALL bracketed placeholders before sending.
%------------------------

\documentclass[a4paper,10pt]{article}

\usepackage{latexsym}
\usepackage[empty]{fullpage}
\usepackage[usenames,dvipsnames]{xcolor}
\usepackage[hidelinks]{hyperref}
\usepackage{fancyhdr}
\usepackage[english]{babel}
\usepackage[T1]{fontenc}
\usepackage{fontawesome5}
\usepackage{eso-pic}
\usepackage{tikz}
\usepackage{graphicx}
\usepackage{geometry}
\input{glyphtounicode}
\pdfgentounicode=1

\geometry{
  a4paper,
  top=0.45in,
  bottom=0.75in,
  left=0.75in,
  right=0.75in
}

\pagestyle{fancy}
\fancyhf{}
\fancyfoot{}
\renewcommand{\headrulewidth}{0pt}
\renewcommand{\footrulewidth}{0pt}

\urlstyle{same}

% ═══════════════════════════════════════════════════════════════════════
% COMPANY DESIGN LAYER
%   cprimary  = header band gradient, left anchor  (dominant brand color)
%   csection  = thin accent rule at band base + achievement label color
%   Band text: always white (#FFFFFF) -- do not change for any company
%   Name text: #FFFFFF  |  Contact text: #FFFFFF  |  Body text: black
%
% WCAG rule of thumb for small contact text (~9pt):
%   R+G+B sum of cprimary hex < 450 -> white text is broadly safe.
%   Examples that pass: #4285F4 (Google Blue, ~4.2:1), #0061FF (Jira),
%   #1E88E5 (Material Blue), #006699 (any mid-navy). If borderline, darken
%   cprimary by 10-15% (add !85!black in the shade definition below).
%
% ┌── SWAP COMPANY: Change only these two hex codes ────────────────────┐
\definecolor{cprimary}{HTML}{4285F4}   % Header gradient anchor (left)
\definecolor{csection}{HTML}{EA4335}   % Accent rule + achievement labels
% └───────────────────────────────────────────────────────────────────┘
% ═══════════════════════════════════════════════════════════════════════

% HEADER BAND (drawn as page background, appears on every page of this letter)
%
% Horizontal gradient:
%   Left  = cprimary (full brand color, strongest brand signal)
%   Right = cprimary!55!black (brand color darkened ~45%)
%
% Why darken the right side?
%   (1) Creates depth and a premium feel -- reads as "designed", not a solid bar.
%   (2) Improves WCAG contrast at the right margin where lighter primaries
%       would otherwise fall below 3:1 against white.
%   (3) The gradient direction (L->R) mirrors natural reading flow.
%
% csection rule (2.5pt) at band base:
%   Visually closes the band and introduces the secondary brand color.
%   This is the same accent color used for achievement labels in the body,
%   creating a consistent two-color brand identity across the document.
\AddToShipoutPictureBG{%
  \begin{tikzpicture}[remember picture, overlay]
    % Gradient band: full page width, 1.55in tall from top
    \shade[left color=cprimary, right color=cprimary!55!black]
      (current page.north west)
      rectangle
      ([yshift=-1.55in]current page.north east);
    % Accent rule at base of band (csection color, 2.5pt)
    \draw[color=csection, line width=2.5pt]
      ([yshift=-1.55in]current page.north west)
      --
      ([yshift=-1.55in]current page.north east);
  \end{tikzpicture}%
}

\begin{document}

% ─────────────────────────────────────────────────────────────────────
% HEADER BLOCK (centered inside the gradient band)
%
% Geometry: top=0.45in places the text area inside the band from line 1.
% Name block (~0.55in tall) + \vspace{0.65in} = ~1.65in total from top,
% clearing the 1.55in band and the 2.5pt rule with breathing room.
% ─────────────────────────────────────────────────────────────────────
\begin{center}
  {\color{white}\textbf{\fontsize{26}{30}\selectfont\scshape Satvik Jain}} \\[5pt]
  {\color{white}\small
    \href{tel:+917678296693}{\faPhone\enspace +91-7678296693}
    \enspace $|$ \enspace
    \href{mailto:satvik.jain@iitdalumni.com}{\faEnvelope\enspace satvik.jain@iitdalumni.com}
    \enspace $|$ \enspace
    \href{https://linkedin.com/in/satvik-jain-iitd}{\faLinkedin\enspace linkedin.com/in/satvik-jain-iitd}
  }
\end{center}

\vspace{0.65in}   % clears band + rule; letter body begins here

% ─────────────────────────────────────────────────────────────────────
% LETTER BODY
% ─────────────────────────────────────────────────────────────────────

\noindent
[DATE -- e.g., 24 May 2026]

\vspace{0.2in}

\noindent
Dear [HIRING MANAGER NAME or Hiring Team],

\vspace{0.2in}

% SECTION 1: HOOK + IDENTITY
%
% Do NOT open with any of these:
%   "I am writing to express my interest in..."
%   "I am excited to apply for..."
%   "I believe I would be a great fit..."
%   "I am passionate about..."
%
% DO open with something specific and real -- a product the company ships,
% a challenge they have named publicly, or a decision they made that you
% noticed. Then introduce yourself through what you build, not a job title.
% Be direct. Two or three sentences is enough.
\noindent
[OPENING PARAGRAPH: 2-3 sentences. Start with a specific named thing about
this company or role. Introduce yourself through your work, not your title.
Human, direct, and impossible to confuse with a template.]

\vspace{0.2in}

% SECTION 2: EVIDENCE (3 achievements, each labeled in brand color)
%
% Label = short noun phrase describing what the achievement is about.
% Body = 2-3 sentences: what you did, what resulted, why it is relevant here.
% One strong metric beats three weak ones. Do not stack metrics for its own sake.
% Write the story of what the problem actually was, not just what the number was.
\noindent
[CONTEXT SENTENCE: one sentence that bridges your background to their specific need.]

\vspace{0.12in}

\noindent
{\textcolor{csection}{\textbf{[ACHIEVEMENT 1 LABEL].}}}
[What you did, what resulted, and why it maps directly to what this role needs.
2-3 sentences. One strong metric. Let the problem context do the heavy lifting.]

\vspace{0.12in}

\noindent
{\textcolor{csection}{\textbf{[ACHIEVEMENT 2 LABEL].}}}
[Stay specific and human. Do not let metrics replace the story of what the problem
actually was and how you decided to approach it the way you did.]

\vspace{0.12in}

\noindent
{\textcolor{csection}{\textbf{[ACHIEVEMENT 3 LABEL].}}}
[End with something that shows how you operate in ambiguity, make ownership calls,
or build things that did not exist before you showed up.]

\vspace{0.2in}

% SECTION 3: COMPANY ALIGNMENT + MOTIVATION
%
% This paragraph must be impossible to copy-paste unchanged to a different
% company's cover letter. Name a product, a stated roadmap direction, or a
% specific challenge they have talked about. Say what draws you there beyond
% the job title and the pay. If you cannot write something specific here,
% you are not ready to write this paragraph yet -- go read more first.
\noindent
[MOTIVATION PARAGRAPH: 2-3 sentences. Something specific about this company.
Why you want to be here, not just why you are qualified to be here. Different companies
build the same products differently -- say which part of how they do it you care about.]

\vspace{0.2in}

% SECTION 4: CLOSE
%
% One sentence. Confident, not desperate. A specific human ask to continue
% the conversation. Do NOT write "thank you for your time and consideration."
% Do NOT write "I look forward to hearing from you." Both read as filler.
% Write something that sounds like you, not like a template.
\noindent
[CLOSING SENTENCE: A confident, specific, human ask to connect.]

\vspace{0.3in}

\noindent
Sincerely,

\vspace{0.35in}

% Optional: signature image -- place sig.png in same folder as .tex file,
% then uncomment the line below.
% \includegraphics[width=1.3in]{sig.png}\\[-4pt]

\noindent
\textbf{Satvik Jain}

\end{document}
